% =====================================================================
%  Platformă de simulare autonomă a jocului Catan cu agenți AI
%  Raport MDS — Anul 2, Semestrul 2 — FMI, Universitatea din București
%
%  COMPILARE: pdfLaTeX (implicit pe Overleaf).
%  Dacă vrei fonturi mai frumoase, schimbă compilatorul în LuaLaTeX
%  din Overleaf (Menu > Compiler) și vezi nota din preambul.
% =====================================================================
\documentclass[11pt,a4paper]{article}

% ---------- Limbă și codare ----------
% Documentul compilează pe ORICE motor. Detectăm automat care rulează.
%
% RECOMANDARE: pune Overleaf pe LuaLaTeX (Menu > Compiler > LuaLaTeX).
% Cu pdfLaTeX, literele ș și ț sunt *desenate* corect, dar sunt scrise în
% stratul de text al PDF-ului ca "s" + virgulă separată. Efect: copy-paste
% și căutarea în PDF returnează "s,i" în loc de "și". LuaLaTeX folosește
% glifele Unicode native și elimină problema.
\usepackage{iftex}
\ifPDFTeX
  \usepackage[utf8]{inputenc}
  \usepackage[T1]{fontenc}
  \usepackage{lmodern}
\else
  \usepackage{fontspec}
  \setmainfont{Latin Modern Roman}
  \setsansfont{Latin Modern Sans}
  \setmonofont{Latin Modern Mono}
\fi
\usepackage[romanian]{babel}

% ---------- Layout ----------
\usepackage[a4paper,margin=2.4cm,headheight=14pt]{geometry}
\usepackage{fancyhdr}
\usepackage{titlesec}
\usepackage{parskip}
\usepackage{microtype}

% ---------- Grafică și culoare ----------
\usepackage{graphicx}
\usepackage{xcolor}
\usepackage{float}
\usepackage[most]{tcolorbox}

% ---------- Tabele ----------
\usepackage{booktabs}
\usepackage{tabularx}
\usepackage{longtable}
\usepackage{array}

% ---------- Cod sursă ----------
\usepackage{listings}

% ---------- Linkuri ----------
\usepackage[hidelinks]{hyperref}

% =====================================================================
%  PALETĂ DE CULORI — folosește ACELEAȘI valori și în draw.io
%  ca diagramele să arate ca parte din același document.
% =====================================================================
\definecolor{bbNavy}{HTML}{0F2942}    % titluri, contur
\definecolor{bbBlue}{HTML}{2563EB}    % accent primar
\definecolor{bbAmber}{HTML}{D97706}   % accent secundar
\definecolor{bbGreen}{HTML}{059669}   % succes / test trecut
\definecolor{bbRed}{HTML}{DC2626}     % bug / eroare
\definecolor{bbGray}{HTML}{64748B}    % text secundar
\definecolor{bbLight}{HTML}{F1F5F9}   % fundal blocuri

% ---------- Stil titluri ----------
\titleformat{\section}
  {\color{bbNavy}\Large\bfseries}{\thesection}{0.6em}{}
  [\vspace{-0.6em}\color{bbBlue}\rule{\textwidth}{1pt}]
\titleformat{\subsection}
  {\color{bbNavy}\large\bfseries}{\thesubsection}{0.6em}{}
\titleformat{\subsubsection}
  {\color{bbGray}\normalsize\bfseries}{\thesubsubsection}{0.6em}{}

% ---------- Antet / subsol ----------
\pagestyle{fancy}
\fancyhf{}
\fancyhead[L]{\small\color{bbGray}Platformă Simulare Autonomă Catan}
\fancyhead[R]{\small\color{bbGray}MDS 2026}
\fancyfoot[C]{\small\color{bbGray}\thepage}
\renewcommand{\headrulewidth}{0.4pt}

% ---------- Stil cod ----------
\lstdefinestyle{bb}{
  basicstyle=\ttfamily\footnotesize,
  keywordstyle=\color{bbBlue}\bfseries,
  commentstyle=\color{bbGray}\itshape,
  stringstyle=\color{bbGreen},
  numberstyle=\tiny\color{bbGray},
  numbers=left, numbersep=8pt,
  backgroundcolor=\color{bbLight},
  frame=leftline, framerule=2pt, rulecolor=\color{bbBlue},
  breaklines=true, breakatwhitespace=false,
  showstringspaces=false,
  tabsize=2, xleftmargin=12pt, framexleftmargin=12pt,
  captionpos=b,
  extendedchars=true,
}
% Sub pdfLaTeX, listings nu digeră UTF-8 în corpul verbatim; îi dăm o hartă
% explicită. Sub LuaLaTeX nu e nevoie — Unicode funcționează nativ.
\ifPDFTeX
\lstset{
  literate=%
    {ă}{{\u{a}}}1 {Ă}{{\u{A}}}1
    {â}{{\^a}}1   {Â}{{\^A}}1
    {î}{{\^\i}}1 {Î}{{\^I}}1
    {ș}{{\c{s}}}1 {Ș}{{\c{S}}}1
    {ț}{{\c{t}}}1 {Ț}{{\c{T}}}1
    {—}{{---}}1    {–}{{--}}1,
}
\fi
\lstset{style=bb}

% ---------- Cutii reutilizabile ----------
\newtcolorbox{userstory}[1]{
  enhanced, breakable,
  colback=bbLight, colframe=bbBlue,
  boxrule=0pt, leftrule=3pt, arc=2pt,
  title={\bfseries #1}, coltitle=bbNavy,
  colbacktitle=bbLight, fonttitle=\normalsize,
  attach boxed title to top left={xshift=6pt,yshift=-2pt},
  top=10pt,
}
\newtcolorbox{infobox}[1]{
  enhanced, breakable,
  colback=white, colframe=bbAmber,
  boxrule=1pt, arc=3pt,
  title={\bfseries #1}, coltitle=white, colbacktitle=bbAmber,
}
\newtcolorbox{bugbox}[1]{
  enhanced, breakable,
  colback=white, colframe=bbRed,
  boxrule=1pt, arc=3pt,
  title={\bfseries #1}, coltitle=white, colbacktitle=bbRed,
}

% Comandă scurtă pentru criterii de acceptare
\newcommand{\ca}[1]{\par\smallskip{\footnotesize\textbf{\color{bbGray}Criterii de acceptare:} #1}}

% Comandă pentru inserarea unei diagrame draw.io exportate ca PDF.
% Dacă fișierul nu există încă, afișează un placeholder în locul lui,
% astfel încât documentul să compileze chiar înainte de a desena diagramele.
\newcommand{\diagrama}[3]{%
  \begin{figure}[H]
    \centering
    \IfFileExists{diagrame/#1.pdf}
      {\includegraphics[width=#2\textwidth,height=0.80\textheight,keepaspectratio]{diagrame/#1}}
      {\IfFileExists{diagrame/#1.png}
        {\includegraphics[width=#2\textwidth,height=0.80\textheight,keepaspectratio]{diagrame/#1}}
        {\begin{tcolorbox}[enhanced,colback=bbLight,colframe=bbGray,
            boxrule=0pt,arc=3pt,width=#2\textwidth,halign=center,
            borderline={1.2pt}{0pt}{bbGray,dashed}]
          \vspace{1.2cm}
          {\color{bbGray}\bfseries\large [ DIAGRAMA LIPSĂ ]}\\[0.4cm]
          {\color{bbGray}\ttfamily diagrame/#1.pdf}\\[0.3cm]
          {\color{bbGray}\small Exportă diagrama din draw.io și încarc-o aici.}
          \vspace{1.2cm}
        \end{tcolorbox}}}
    \caption{#3}
    \label{fig:#1}
  \end{figure}%
}

\begin{document}

% =====================================================================
%  PAGINA DE TITLU
% =====================================================================
\begin{titlepage}
\centering
\vspace*{2cm}

{\color{bbBlue}\rule{\textwidth}{2pt}}
\vspace{0.8cm}

{\Huge\bfseries\color{bbNavy}
Platformă de simulare autonomă\\[0.3cm]
a jocului de societate Catan\\[0.3cm]
cu agenți AI\par}

\vspace{0.6cm}
{\color{bbBlue}\rule{\textwidth}{2pt}}

\vspace{1.5cm}
{\Large\color{bbGray}Raport de proiect\par}
\vspace{0.3cm}
{\large\color{bbGray}Metode de Dezvoltare Software\par}

\vspace{2.5cm}

\begin{tabular}{r@{\hspace{1em}}l}
\textbf{\color{bbNavy}Echipă:} & Rusu Iulia \\
                               & Barcan Silviu-Ioan \\
                               & Cocut Ioana \\[0.6cm]
\textbf{\color{bbNavy}Disciplina:} & MDS, Anul 2, Semestrul 2 \\[0.2cm]
\textbf{\color{bbNavy}Facultatea:} & Matematică și Informatică, \\
                                   & Universitatea din București \\[0.2cm]
\textbf{\color{bbNavy}Data:} & Septembrie 2026 \\
\end{tabular}

\vfill

\begin{tcolorbox}[colback=bbLight,colframe=bbNavy,boxrule=1pt,arc=3pt,width=0.85\textwidth]
\centering\small
\textbf{Repository:} \url{https://github.com/sigutz/Byte-Boards}\\[0.2cm]
\textbf{Producție:} \url{https://www.bytenboards.games}
\end{tcolorbox}

\vspace{1cm}
\end{titlepage}

\tableofcontents
\newpage

% =====================================================================
\section{Introducere și context}
% =====================================================================

\subsection{Descrierea produsului}

Proiectul implementează o platformă web în care agenți AI joacă autonom
\textit{Catan}, iar utilizatorii urmăresc partidele în timp real. Spre deosebire de o
implementare clasică de joc, utilizatorul nu joacă --- el configurează agenții,
pornește meciul și observă deciziile luate de aceștia, însoțite de comentariu
generat automat.

Sistemul suportă două seturi de reguli: \textbf{Catan Classic} (tablă de 19 hexagoane,
54 de noduri) și \textbf{Catan Seafarers} (extindere cu 3 insule exterioare,
nave, pirat și puncte de victorie pentru explorare).

\subsection{Stack tehnologic}

\begin{table}[H]
\centering
\begin{tabularx}{\textwidth}{@{}l l X@{}}
\toprule
\textbf{Nivel} & \textbf{Tehnologie} & \textbf{Rol în sistem} \\
\midrule
Frontend & React 19 + Vite 7 & SPA cu routing manual (\texttt{pushState}), fără librărie de router \\
Backend  & Express 5 + Node.js & API REST (24 de rute) și flux SSE \\
ORM      & Prisma 7 + \texttt{@prisma/adapter-pg} & 9 modele; pool \texttt{pg} partajat pentru SQL brut \\
Bază de date & PostgreSQL & Persistență meciuri, agenți, evenimente \\
AI        & Google Gemini & Decizii agenți, comentariu, rezumate, moderare \\
Infrastructură & Docker + Nginx & Build multi-stage, reverse proxy \\
CI/CD     & GitHub Actions & 6 job-uri de verificare + deploy prin SSH \\
\bottomrule
\end{tabularx}
\caption{Stack-ul tehnologic al platformei}
\end{table}

\subsection{Metodologia de dezvoltare}

Echipa a folosit un proces \textbf{iterativ-incremental} de tip Agile, cu
livrare pe funcționalități:

\begin{itemize}
  \item cerințele au fost formulate ca \textbf{user stories} cu criterii de
        acceptare explicite, grupate pe \textbf{epice};
  \item prioritizarea backlog-ului s-a făcut prin metoda \textbf{MoSCoW};
  \item fiecare funcționalitate a fost dezvoltată pe un \textbf{feature branch}
        izolat și integrată în \texttt{main} printr-un \textbf{Pull Request} cu review;
  \item integrarea a fost protejată de un \textbf{pipeline CI} care rulează
        teste, build și lint la fiecare push.
\end{itemize}

\newpage
% =====================================================================
\section{User Stories și Backlog}
% =====================================================================

Backlog-ul a fost organizat pe \textbf{4 epice}, cu 15 user stories. Fiecare
poveste identifică explicit \textbf{actorul} (Vizitator, Utilizator autentificat,
Administrator), ceea ce a permis derivarea directă a diagramei de cazuri de
utilizare din Secțiunea~\ref{sec:usecase}.

\subsection{Epic 1 --- Configurarea și urmărirea live a meciurilor}

\begin{userstory}{US-01 --- Urmărirea unui meci live}
\textbf{Ca} Vizitator cu link către un meci, \textbf{vreau} să urmăresc agenții AI
jucând în timp real, \textbf{pentru ca} să pot observa deciziile și strategiile lor.
\ca{Tabla se actualizează automat, fără refresh manual; fiecare mutare apare în
maximum 3 secunde; fără link, un Vizitator nu poate descoperi meciuri.}
\end{userstory}

\begin{userstory}{US-02 --- Alegerea jocului și a modului}
\textbf{Ca} Utilizator autentificat, \textbf{vreau} să aleg între Catan Classic și
Catan Seafarers, \textbf{pentru ca} să pot selecta setul de reguli dorit.
\ca{Selector vizibil la crearea meciului; în Seafarers sunt disponibile nave și
puncte de victorie pentru insule.}
\end{userstory}

\begin{userstory}{US-03 --- Pornirea unui meci la cerere}
\textbf{Ca} Utilizator autentificat, \textbf{vreau} să pornesc un meci nou alegând
agenții participanți, \textbf{pentru ca} să generez conținut nou de urmărit.
\ca{Minimum 2 agenți eligibili; limita de agenți este impusă în UI cu mesaj clar.}
\end{userstory}

\begin{userstory}{US-04 --- Oprirea sau ștergerea unui meci}
\textbf{Ca} Utilizator autentificat, \textbf{vreau} să opresc sau să șterg un meci
pe care îl dețin, \textbf{pentru ca} să pot anula un meci blocat sau pornit din greșeală.
\ca{La „Stop” meciul se oprește și e marcat încheiat; la ștergere dispare din istoric.}
\end{userstory}

\begin{userstory}{US-05 --- Comentariu live generat de AI}
\textbf{Ca} Vizitator, \textbf{vreau} să văd comentariu generat de AI în timpul
meciului, \textbf{pentru ca} experiența să fie mai captivantă.
\ca{La un eveniment semnificativ (comerț, hoț) apare o linie de comentariu în
câteva secunde.}
\end{userstory}

\begin{userstory}{US-06 --- Vizualizarea trăsăturilor agenților}
\textbf{Ca} Vizitator, \textbf{vreau} să văd trăsăturile fiecărui agent,
\textbf{pentru ca} să înțeleg \textit{de ce} ia o anumită decizie, nu doar \textit{ce} face.
\ca{Trăsăturile sunt afișate per agent; trăsăturile în conflict sunt rezolvate
conform regulilor \texttt{TraitConflict} fără eroare.}
\end{userstory}

\subsection{Epic 2 --- Istoric și analitice}

\begin{userstory}{US-07 --- Navigarea în istoricul meciurilor}
\textbf{Ca} Utilizator autentificat, \textbf{vreau} să răsfoiesc meciurile trecute,
\textbf{pentru ca} să pot urmări tipare în timp.
\ca{Listă cu dată, tip de joc și agenți participanți.}
\end{userstory}

\begin{userstory}{US-08 --- Detalii și statistici de meci}
\textbf{Ca} Utilizator autentificat, \textbf{vreau} să deschid un meci trecut,
\textbf{pentru ca} să analizez momentele cheie.
\ca{Log complet de evenimente per jucător, deținătorul drumului cel mai lung,
scoruri finale.}
\end{userstory}

\begin{userstory}{US-09 --- Filtrarea istoricului}
\textbf{Ca} Utilizator autentificat, \textbf{vreau} să filtrez istoricul după tipul
de joc, \textbf{pentru ca} să mă concentrez pe ce mă interesează.
\ca{La selectarea unui filtru sunt afișate doar meciurile de acel tip.}
\end{userstory}

\begin{userstory}{US-10 --- Metrici de performanță per agent}
\textbf{Ca} Utilizator autentificat, \textbf{vreau} să văd statistici agregate per
agent, \textbf{pentru ca} să înțeleg care performează cel mai bine.
\ca{Rată de câștig, număr de meciuri, scor mediu, defalcate pe tip de joc.}
\end{userstory}

\subsection{Epic 3 --- Autentificare și proprietate}

\begin{userstory}{US-11 --- Înregistrare și autentificare}
\textbf{Ca} Vizitator, \textbf{vreau} să îmi creez cont și să mă autentific,
\textbf{pentru ca} să dețin agenți și să controlez cine îmi vede meciurile.
\ca{Parola stocată cu \texttt{bcrypt}; sesiune JWT; la credențiale greșite nu se
creează sesiune.}
\end{userstory}

\begin{userstory}{US-12 --- Deținerea agenților proprii}
\textbf{Ca} Utilizator autentificat, \textbf{vreau} să dețin agenții pe care îi
creez, \textbf{pentru ca} doar eu să îi pot configura sau înscrie în meciuri.
\ca{Agenții de sistem sunt vizibili tuturor; agenții proprii doar creatorului;
numele sunt moderate automat înainte de salvare.}
\end{userstory}

\subsection{Epic 4 --- Confidențialitate, partajare și administrare}

\begin{userstory}{US-13 --- Controlul vizibilității meciurilor}
\textbf{Ca} Utilizator autentificat, \textbf{vreau} să marchez meciurile ca private,
protejate sau publice, \textbf{pentru ca} eu să decid cine le vede.
\ca{Un utilizator fără drept de acces nu poate vedea conținutul unui meci privat.}
\end{userstory}

\begin{userstory}{US-14 --- Partajare prin link sau invitație}
\textbf{Ca} Utilizator autentificat, \textbf{vreau} să generez un link de partajare
bazat pe token, \textbf{pentru ca} să pot arăta un meci interesant unor persoane anume.
\ca{Cu tokenul valid, meciul poate fi vizualizat fără autentificare.}
\end{userstory}

\begin{userstory}{US-15 --- Moderare de către administrator}
\textbf{Ca} Administrator, \textbf{vreau} să pot șterge orice meci indiferent de
proprietar, \textbf{pentru ca} să elimin conținut nepotrivit sau defect.
\ca{Rolul \texttt{ADMIN} are drept de ștergere pe orice resursă.}
\end{userstory}

\subsection{Prioritizare prin metoda MoSCoW}

\begin{table}[H]
\centering
\begin{tabularx}{\textwidth}{@{}>{\bfseries}l X@{}}
\toprule
\textbf{Prioritate} & \textbf{User Stories} \\
\midrule
\textcolor{bbRed}{Must have}    & US-01, US-02, US-03, US-05, US-07, US-08, US-11 \\
\textcolor{bbAmber}{Should have} & US-04, US-09, US-10, US-12 \\
\textcolor{bbBlue}{Could have}  & US-06, US-13, US-14, US-15 \\
\textcolor{bbGray}{Won't have (v1)} & Multiplayer uman, chat live între utilizatori, turnee \\
\bottomrule
\end{tabularx}
\caption{Prioritizarea cerințelor prin metoda MoSCoW}
\end{table}

\newpage
% =====================================================================
\section{Diagrame UML}
% =====================================================================

Modelarea sistemului folosește notația \textbf{UML 2.5}. Diagramele acoperă
atât perspectiva \textbf{structurală} (cazuri de utilizare, clase, componente,
desfășurare), cât și cea \textbf{comportamentală} (secvență, stare, activitate).
Toate au fost realizate în \texttt{draw.io} și exportate vectorial (PDF), astfel
încât să rămână clare la orice nivel de zoom.

% ---------------------------------------------------------------------
\subsection{Diagrama cazurilor de utilizare}
\label{sec:usecase}

Diagrama de cazuri de utilizare formalizează cerințele din Secțiunea 2. Sistemul
are trei actori umani cu drepturi cumulative --- \textbf{Vizitator}
$\subset$ \textbf{Utilizator autentificat} $\subset$ \textbf{Administrator} ---
și un \textbf{actor secundar non-uman}, API-ul Gemini, invocat de sistem pentru
deciziile agenților și generarea de text.

\diagrama{01-use-case}{1.0}{Diagrama cazurilor de utilizare --- actori și funcționalități}

Relațiile modelate sunt:
\begin{itemize}
  \item \textbf{Generalizare} între actori: Administratorul moștenește tot ce
        poate face un Utilizator autentificat, care la rândul lui moștenește
        drepturile Vizitatorului.
  \item \texttt{\textless\textless include\textgreater\textgreater} ---
        \textit{Pornire meci} include obligatoriu \textit{Selectare agenți}
        și \textit{Moderare nume}.
  \item \texttt{\textless\textless extend\textgreater\textgreater} ---
        \textit{Generare rezumat AI} extinde \textit{Finalizare meci}, fiind
        opțională (la indisponibilitatea API-ului se folosește un text implicit).
\end{itemize}

% ---------------------------------------------------------------------
\subsection{Diagrama de clase (modelul de domeniu)}

Schema relațională este descrisă de 9 entități. Relația centrală este
\texttt{MatchAgent} --- o \textbf{clasă de asociere} între \texttt{Match} și
\texttt{Agent}, care nu doar leagă cele două entități, ci poartă întreaga stare
de joc a unui jucător într-un meci (resurse, clădiri, cărți de dezvoltare).

\diagrama{02-class-diagram}{1.0}{Diagrama de clase --- modelul de domeniu cu multiplicități}

\begin{infobox}{Decizie de modelare}
Trăsăturile agenților sunt modelate \textbf{relațional}, nu ca un câmp JSON:
\texttt{Trait} (catalogul), \texttt{AgentTrait} (asocierea many-to-many) și
\texttt{TraitConflict} (auto-asociere pe \texttt{Trait} care codifică perechile
incompatibile). Astfel validarea combinațiilor invalide --- de exemplu
\textit{Empathic} împreună cu \textit{Aggressive} --- se face prin interogare,
nu prin logică hardcodată în aplicație.
\end{infobox}

% ---------------------------------------------------------------------
\subsection{Diagrama de componente}

Diagrama de componente arată descompunerea logică a sistemului și
\textbf{interfețele oferite și cerute} între module.

\diagrama{03-component}{1.0}{Diagrama de componente --- module și interfețe}

Motorul de joc este izolat complet de persistență: modulele din
\texttt{backend/src/game/} sunt funcții pure, fără dependențe de Prisma sau
Express. Această separare este ceea ce permite rularea unui meci complet în
teste, fără bază de date (vezi Secțiunea~\ref{sec:teste}).

% ---------------------------------------------------------------------
\subsection{Diagrama de desfășurare}

Diagrama de desfășurare descrie topologia reală din producție: noduri fizice,
containere Docker, rețele și protocoale.

\diagrama{04-deployment}{1.0}{Diagrama de desfășurare --- containere, rețele și porturi}

\begin{infobox}{Detaliu critic de infrastructură}
Containerul \texttt{backend} este atașat la \textbf{două} rețele Docker:
\texttt{app-net} (comunicarea cu frontend-ul) și \texttt{global-db-net} (rețea
externă, partajată cu containerul de bază de date \texttt{universal-db}).
Omiterea primei rețele a cauzat o cădere completă de producție, analizată în
Secțiunea~\ref{sec:bug3}.
\end{infobox}

% ---------------------------------------------------------------------
\subsection{Diagrama de secvență --- ciclul de viață al unui meci}

Diagrama de secvență urmărește un meci de la cererea HTTP inițială până la
persistarea rezumatului final, evidențiind natura \textbf{asincronă} a simulării:
răspunsul HTTP se întoarce imediat, iar simularea continuă în fundal, comunicând
cu clientul prin Server-Sent Events.

\diagrama{05-sequence}{0.95}{Diagrama de secvență --- pornirea și derularea unui meci}

Elementele importante sunt fragmentele combinate:
\begin{itemize}
  \item \texttt{loop} --- bucla de simulare, repetată până la atingerea a 8 puncte
        de victorie sau a limitei de 500 de ture;
  \item \texttt{alt} --- ramificarea pe disponibilitatea API-ului Gemini: dacă
        apelul eșuează sau lipsește cheia, se aplică euristica deterministă;
  \item mesajul \textbf{asincron} (săgeată deschisă) pentru evenimentele SSE,
        spre deosebire de apelurile sincrone (săgeată plină).
\end{itemize}

% ---------------------------------------------------------------------
\subsection{Diagrama de stare --- statusul unui meci}

Entitatea \texttt{Match} are un ciclu de viață explicit, codificat în enumerarea
\texttt{MatchStatus}.

\diagrama{06-state-machine}{0.70}{Diagrama de stare --- tranzițiile unui meci}

Tranziția \texttt{LIVE} $\rightarrow$ \texttt{PAUSED} $\rightarrow$ \texttt{LIVE}
este reversibilă, în timp ce \texttt{COMPLETED} este o \textbf{stare finală}:
odată atinsă, meciul nu mai poate fi repornit. Repornirea serverului declanșează
tranziția de recuperare --- meciurile rămase \texttt{LIVE} în baza de date sunt
reluate automat din starea persistată.

% ---------------------------------------------------------------------
\subsection{Diagrama de activitate --- algoritmul unei ture}

O tură se desfășoară în două faze distincte: mai întâi \textbf{producția}, comună
tuturor jucătorilor și determinată de zaruri, apoi \textbf{decizia}, repetată
individual pentru fiecare agent. Cele două sunt prezentate separat, pentru
lizibilitate.

\subsubsection{Faza de producție}

\diagrama{07a-activity-productie}{0.55}{Diagrama de activitate --- faza de producție a resurselor}

Nodul de decizie pe rezultatul zarurilor ($=7$ sau $\neq 7$) separă cele două
regimuri de joc: producția normală de resurse, respectiv activarea hoțului cu
descărcarea obligatorie a jucătorilor care dețin cel puțin 7 resurse. Această
fază se execută o singură dată per tură, înaintea oricărei decizii de agent.

\subsubsection{Faza de decizie a agenților}

\diagrama{07b-activity-decizie}{0.50}{Diagrama de activitate --- decizia și acțiunea unui agent}

Bucla parcurge agenții în ordinea locurilor. Pentru fiecare, motorul calculează
mai întâi mulțimea mutărilor legale; abia apoi intervine componenta AI, iar
nodul de decizie pe disponibilitatea cheii Gemini arată explicit calea de
rezervă: la eșec sau lipsă de cheie se aplică euristica deterministă, astfel
încât simularea nu se blochează niciodată.

După aplicarea acțiunii, persistarea stării în baza de date și emiterea
evenimentului SSE sunt \textbf{independente} între ele --- niciuna nu așteaptă
rezultatul celeilalte. Ieșirea din buclă se face fie la atingerea pragului de
victorie, fie la depășirea limitei de ture.

\newpage
% =====================================================================
\section{Source Control cu Git}
% =====================================================================

\subsection{Strategia de branching}

Proiectul folosește un model de tip \textbf{feature branching}: nicio modificare
nu se face direct pe \texttt{main}, iar fiecare funcționalitate majoră este
dezvoltată izolat și integrată printr-un Pull Request cu review.

\diagrama{08-git-flow}{1.0}{Strategia de branching --- ramuri de funcționalitate și puncte de integrare}

\begin{table}[H]
\centering
\begin{tabularx}{\textwidth}{@{}l X l@{}}
\toprule
\textbf{Branch} & \textbf{Scop} & \textbf{Autor} \\
\midrule
\texttt{branch-initial-setup} & Schelet inițial frontend + backend & Rusu Iulia \\
\texttt{login-and-share} & Autentificare JWT, share links, Docker & Rusu Iulia \\
\texttt{ai-branch} & Motor AI Gemini, graf tablă, personalitate & Barcan Silviu \\
\texttt{testeAgenti} & Suită de teste pentru logica jocului & Cocut Ioana \\
\texttt{CItests} & Workflow GitHub Actions & Cocut Ioana \\
\texttt{bot-ownership-admin-controls} & Proprietate agenți, rol admin, privacy & Rusu Iulia \\
\texttt{seafarers-game-mode} & Modul Seafarers complet cu AI agentiv & Rusu Iulia \\
\texttt{tests-and-evals} & Suită Vitest + evals DeepEval & Rusu Iulia \\
\texttt{fix/production-deploy} & Remediere cădere de producție & Barcan Silviu \\
\bottomrule
\end{tabularx}
\caption{Ramurile create în cadrul proiectului}
\end{table}

\subsection{Pull Requests}

Au fost integrate \textbf{13 Pull Requests}, toate cu review înainte de merge.

\begin{table}[H]
\centering
\begin{tabularx}{\textwidth}{@{}l l X@{}}
\toprule
\textbf{PR} & \textbf{Branch} & \textbf{Descriere} \\
\midrule
\#2  & \texttt{branch-initial-setup} & Setup inițial de proiect \\
\#7  & \texttt{login-and-share} & Autentificare, Docker, partajare \\
\#8, \#9 & \texttt{ai-branch} & AI Gemini, graf tablă, personalitate \\
\#10 & \texttt{testeAgenti} & Teste automate pentru agenți \\
\#11 & \texttt{CItests} & Pipeline CI GitHub Actions \\
\#12--\#14 & \texttt{bot-ownership-admin-controls} & Proprietate agenți + administrare \\
\#15, \#16 & \texttt{seafarers-game-mode} & Modul Seafarers \\
\#17 & \texttt{tests-and-evals} & 140 teste Vitest + 3 suite de evals \\
\bottomrule
\end{tabularx}
\caption{Pull Requests integrate în \texttt{main}}
\end{table}

\subsection{Tehnici Git utilizate}

\begin{itemize}
  \item \textbf{Feature branches} --- nicio modificare direct pe \texttt{main};
  \item \textbf{Pull Requests cu code review} --- toate integrările trec prin PR;
  \item \textbf{Rebase} --- sincronizare cu \texttt{main} înainte de PR
        (\texttt{git pull --rebase origin main});
  \item \textbf{Merge commits} --- păstrate în istoric pentru trasabilitate;
  \item \textbf{Branch protection prin CI} --- un PR nu poate fi integrat
        dacă vreo verificare automată eșuează.
\end{itemize}

% =====================================================================
\section{Teste automate și evaluări AI}
\label{sec:teste}
% =====================================================================

Testarea are \textbf{trei niveluri}, fiecare cu un scop distinct: teste
unitare pe motorul de joc, teste de integrare pe API, și \textit{evals}
pentru componenta non-deterministă (LLM).

\begin{table}[H]
\centering
\begin{tabularx}{\textwidth}{@{}l l r X@{}}
\toprule
\textbf{Nivel} & \textbf{Unealtă} & \textbf{Teste} & \textbf{Ce validează} \\
\midrule
Unitare & Vitest & 140 & Reguli de joc, graf tablă, trăsături, simulare completă \\
Integrare & Vitest + Supertest & (incluse) & Rute API, autorizare, degradare fără DB \\
Regresie & Node \texttt{assert} & 15 & Suita istorică \texttt{test-agents.ts} \\
Evals AI & DeepEval (Python) & 3 suite & Calitatea deciziilor și textului generat de LLM \\
\bottomrule
\end{tabularx}
\caption{Nivelurile de testare din proiect}
\end{table}

\subsection{Teste unitare pe motorul de joc}

Modulele din \texttt{backend/src/game/} sunt funcții pure, deci pot fi testate
fără bază de date sau server. Suita acoperă validarea acțiunilor legale
(\texttt{getValidActions}), aplicarea efectelor (\texttt{applyAction}),
calculul punctajului (\texttt{computeVP}), construcția grafului tablei
și sistemul de trăsături.

\begin{lstlisting}[language=bash,caption={Rularea suitei complete de teste}]
cd backend && npm test

  v src/game/personality.test.ts       (17 tests)
  v src/game/graph.test.ts             (21 tests)
  v src/game/seafarers-graph.test.ts   (16 tests)
  v src/game/seafarers.test.ts         (18 tests)
  v src/game/catan.test.ts             (49 tests)
  v src/game/game-simulation.test.ts   (10 tests)
  v src/api.test.ts                    ( 9 tests)

  Test Files  7 passed (7)
       Tests  140 passed (140)
\end{lstlisting}

\subsection{Testarea prin simulare de meci complet}

Cel mai valoros test nu verifică o funcție, ci o \textbf{proprietate invariantă}
pe parcursul unui meci întreg. Modulul \texttt{simulation-harness.ts} rulează un
meci complet folosind exact aceleași funcții ca aplicația reală, dar fără
persistență, iar testele verifică invarianți:

\begin{itemize}
  \item există întotdeauna un câștigător (prin scor sau prin limita de ture);
  \item niciun jucător nu ajunge la o cantitate negativă de resurse;
  \item un nod nu poate fi simultan colonie și oraș;
  \item cel mult un jucător deține \textit{Longest Road} la un moment dat,
        și doar cu un drum de lungime $\geq 5$;
  \item cel mult un jucător deține \textit{Largest Army}, și doar cu
        $\geq 3$ cavaleri jucați.
\end{itemize}

Această abordare --- \textit{property-based testing} --- prinde clase întregi de
erori pe care testele punctuale nu le-ar detecta, pentru că explorează stări de
joc generate aleator la fiecare rulare.

\subsection{Evals pentru agenți AI}

Componenta AI este \textbf{non-deterministă}: nu se poate scrie o aserțiune de
egalitate pe ieșirea unui LLM. Soluția standard este evaluarea prin
\textit{metrici de calitate}, unde un al doilea model joacă rol de
\textbf{judecător}. Proiectul folosește framework-ul \textbf{DeepEval} cu metrici
de tip \texttt{GEval}, rulate ca teste \texttt{pytest} împotriva unor rute
dedicate din backend.

\begin{table}[H]
\centering
\begin{tabularx}{\textwidth}{@{}l X l@{}}
\toprule
\textbf{Eval} & \textbf{Ce măsoară} & \textbf{Prag} \\
\midrule
\texttt{TraitAlignment} & Acțiunea aleasă este consistentă cu trăsăturile declarate ale agentului & 0.6 \\
\texttt{CommentaryGrounding} & Comentariul descrie fidel acțiunea, fără detalii inventate & 0.6 \\
\texttt{Moderare nume} & Acuratețea deciziilor de moderare pe 14 cazuri de test & $\geq$ 90\% \\
\bottomrule
\end{tabularx}
\caption{Metricile de evaluare a componentei AI}
\end{table}

\begin{infobox}{Garanție structurală înaintea evaluării calitative}
Independent de calitatea răspunsului LLM, sistemul are o garanție \textbf{tare}:
\texttt{getAgentDecision()} returnează întotdeauna o acțiune din lista de mutări
legale. Dacă modelul propune o acțiune invalidă, dacă răspunsul nu e JSON valid,
sau dacă apelul eșuează, se aplică euristica deterministă. \textbf{Nicio ieșire a
modelului nu poate corupe starea jocului} --- evals-urile măsoară \textit{cât de
bună} este decizia, nu \textit{dacă este legală}.
\end{infobox}

\newpage
% =====================================================================
\section{Raportare de bug și rezolvare prin Pull Request}
% =====================================================================

\subsection{Bug \#1 --- Eșec de parsare JSON la autentificare}

\begin{bugbox}{Simptom}
După un deploy, pagina de login returna\\
\texttt{JSON.parse: unexpected character at line 1 column 1 of the JSON data}.
\end{bugbox}

\textbf{Investigație.} Frontend-ul apela \texttt{POST /api/auth/login} și încerca
să parseze răspunsul ca JSON, dar Nginx returna o pagină HTML de eroare
\texttt{502 Bad Gateway}. Cauza reală era că backend-ul nu pornise: procesul se
oprea la \texttt{prisma db push}, pentru că schema conținea o coloană cu date
existente pe care Prisma refuza să o șteargă fără confirmare explicită.

\textbf{Cauză rădăcină.} Dockerfile-ul de producție nu includea flagul
\texttt{--accept-data-loss}. \textbf{Rezolvare:} PR \#15, cu reproducere locală
și review.

\subsection{Bug \#2 --- Crash de runtime în modul Seafarers}

\begin{bugbox}{Simptom}
Meciurile Seafarers se opreau după câteva ture cu\\
\texttt{TypeError: Cannot read properties of undefined (reading 'adjacentNodes')}.
\end{bugbox}

\textbf{Investigație.} Funcția \texttt{getValidActions} caută nodurile adiacente
în structura \texttt{BOARD\_NODES}, definită pentru tabla clasică (nodurile
0--53). În Seafarers, jucătorii pot coloniza nodurile exterioare 54--71, care nu
existau în acea structură; accesarea returna \texttt{undefined}.

\begin{lstlisting}[language=Java,caption={Fix aplicat --- optional chaining pentru noduri exterioare}]
// INAINTE:
for (const adjNode of BOARD_NODES[nodeId].adjacentNodes) {

// DUPA:
for (const adjNode of (BOARD_NODES[nodeId]?.adjacentNodes ?? [])) {
\end{lstlisting}

Nodurile exterioare sunt astfel ignorate la calculul drumurilor --- ceea ce este
corect din punct de vedere al regulilor, drumurile neputând fi construite pe mare.
\textbf{Rezolvare:} PR \#15.

\subsection{Bug \#3 --- Cădere completă de producție: CI verde, imagine nefuncțională}
\label{sec:bug3}

Acesta este cel mai instructiv defect al proiectului, pentru că ilustrează o
categorie de erori pe care testele unitare \textbf{nu o pot prinde prin
construcție}: o divergență între mediul de test și artefactul livrat.

\begin{bugbox}{Simptom}
Ambele containere în restart infinit. Backend:
\texttt{Cannot find module '/app/dist/src/index.js'}.
Frontend: \texttt{[emerg] host not found in upstream "backend"}.
\end{bugbox}

\subsubsection{Cauză rădăcină A --- calea artefactului depinde de contextul de build}

Fișierul \texttt{tsconfig.json} nu declara explicit \texttt{rootDir}. În lipsa
lui, TypeScript îl \textbf{deduce} din prefixul comun al fișierelor de intrare ---
deci calea de ieșire depinde de \textit{ce fișiere sunt prezente}:

\begin{table}[H]
\centering
\begin{tabularx}{\textwidth}{@{}l X l@{}}
\toprule
\textbf{Context} & \textbf{Fișiere de intrare} & \textbf{Ieșire \texttt{tsc}} \\
\midrule
Checkout complet (CI, local) & \texttt{src/**/*} + \texttt{test-agents.ts} & \texttt{dist/src/index.js} \\
Imagine Docker & doar \texttt{src/**/*} & \texttt{dist/index.js} \\
\bottomrule
\end{tabularx}
\caption{Divergența căii de ieșire între CI și imaginea Docker}
\end{table}

Dockerfile-ul copiază doar directorul \texttt{src/}, nu și \texttt{test-agents.ts}.
Un commit anterior eliminase \texttt{rootDir}; un commit ulterior a „corectat”
comanda de pornire către \texttt{dist/src/index.js}, potrivind-o cu ce producea
CI-ul --- dar nu cu ce producea imaginea. \textbf{Rezultat: CI verde, imagine care
nu putea porni niciodată.}

\subsubsection{Cauză rădăcină B --- servicii pe rețele Docker disjuncte}

În \texttt{docker-compose.yml}, serviciul \texttt{backend} declara doar rețeaua
\texttt{global-db-net}. În Docker Compose, \textbf{declararea oricărei rețele
înlocuiește rețeaua implicită}, deci \texttt{frontend} și \texttt{backend} nu mai
aveau nicio rețea comună, iar numele \texttt{backend} devenea nerezolvabil.
Nginx rezolvă gazda dintr-un \texttt{proxy\_pass} literal \textbf{o singură dată,
la pornire}, și refuză să pornească dacă rezoluția eșuează --- astfel cădea și
frontend-ul.

\subsubsection{Remediere}

\begin{lstlisting}[language=Java,caption={Fixarea căii de ieșire a compilatorului}]
// tsconfig.json -- rootDir fixat explicit
"rootDir": "./src",     // iesirea nu mai depinde de contextul de build
"include": ["src/**/*"],
"exclude": ["node_modules", "dist", "src/**/*.test.ts"]
\end{lstlisting}

\begin{lstlisting}[language=bash,caption={Verificare la build --- imaginea eșuează în loc să pornească defect}]
RUN npm run build
RUN test -f dist/index.js
\end{lstlisting}

Pentru rețea s-a introdus o rețea comună \texttt{app-net}, iar configurația Nginx
a fost transformată în \textit{template} randat de entrypoint-ul imaginii, cu
rezolvare \textbf{leneșă} a upstream-ului printr-o variabilă. Efect: Nginx
pornește chiar și fără backend disponibil, degradează la \texttt{502} și
\textbf{își revine automat} când backend-ul redevine disponibil, fără restart.

\subsubsection{Prevenirea recurenței}

Golul real era că CI valida \texttt{npm run build} pe un checkout complet ---
ceea ce \textbf{nu} este ce face imaginea. S-a adăugat un job care construiește
imaginile reale, verifică existența punctului de intrare și validează
configurația Nginx fără upstream disponibil:

\begin{lstlisting}[language=bash,caption={Job-ul CI care previne recurența}]
- run: docker build -t bytenboard-backend ./backend
- run: docker run --rm bytenboard-backend test -f dist/index.js
- run: docker build --build-arg VITE_API_BASE_URL= -t bytenboard-frontend ./frontend
- run: docker run --rm bytenboard-frontend nginx -t
\end{lstlisting}

\begin{infobox}{Lecție de proces}
Un pipeline CI nu validează codul, ci \textbf{artefactul pe care îl livrează}.
Dacă mediul de test diferă de mediul de producție --- fie și doar prin lista de
fișiere copiate --- testele pot trece în timp ce produsul livrat este nefuncțional.
\end{infobox}

\newpage
% =====================================================================
\section{Pipeline CI/CD}
% =====================================================================

Proiectul are două workflow-uri GitHub Actions distincte: \texttt{test.yml}
(rulează la orice push sau PR, pe orice ramură) și \texttt{deploy.yml}
(rulează doar la push pe \texttt{main} și livrează în producție).

\diagrama{09-cicd}{1.0}{Pipeline-ul CI/CD --- verificări automate și livrare}

\subsection{Verificări automate}

\begin{table}[H]
\centering
\begin{tabularx}{\textwidth}{@{}l X X@{}}
\toprule
\textbf{Job} & \textbf{Ce verifică} & \textbf{Eșuează dacă} \\
\midrule
\texttt{test-agents} & Suita istorică de regresie & O aserțiune eșuează \\
\texttt{test-backend} & 140 de teste Vitest & Orice test pică \\
\texttt{build-backend} & TypeScript compilabil & Erori de tip sau import \\
\texttt{build-frontend} & Bundle Vite compilabil & Erori TS, importuri lipsă \\
\texttt{lint-frontend} & ESLint & Variabile neutilizate, tipuri \texttt{any} \\
\texttt{docker-build} & \textbf{Imaginile reale de producție} & Punctul de intrare lipsește sau Nginx e invalid \\
\midrule
\texttt{deploy} & Livrare prin SSH & Oricare dintre cele de mai sus eșuează \\
\bottomrule
\end{tabularx}
\caption{Verificările automate din pipeline}
\end{table}

Ultimele două job-uri (\texttt{test-backend}, \texttt{docker-build}) au fost
adăugate ca urmare directă a incidentului din Secțiunea~\ref{sec:bug3}.

\subsection{Gestionarea secretelor}

Pipeline-ul injectează secretele din GitHub Secrets într-un fișier \texttt{.env}
generat pe server la fiecare deploy:

\begin{itemize}
  \item \texttt{DATABASE\_URL}, \texttt{JWT\_SECRET}, \texttt{VITE\_API\_BASE\_URL};
  \item \texttt{GEMINI\_1} \dots \texttt{GEMINI\_8} --- pool-ul de 8 chei API;
  \item \texttt{SERVER\_HOST}, \texttt{SERVER\_USER}, \texttt{SERVER\_SSH\_KEY}.
\end{itemize}

Niciun secret nu este versionat: \texttt{.env} figurează în \texttt{.gitignore}.

\subsection{Impactul pipeline-ului}

Politica „niciun merge fără verificări verzi” a prins activ:
\begin{itemize}
  \item erori TypeScript introduse de câmpuri noi (\texttt{shipEdges},
        \texttt{islandVPs}) neadăugate în factory-ul de test;
  \item erori ESLint preexistente, expuse de verificarea adăugată în PR \#12;
  \item blocaje de deploy cauzate de un build de frontend eșuat.
\end{itemize}

% =====================================================================
\section{Folosirea instrumentelor AI în dezvoltare}
% =====================================================================

Proiectul implică AI pe \textbf{două axe distincte}, care nu trebuie confundate:
AI ca \textit{instrument de dezvoltare} și AI ca \textit{funcționalitate a
produsului}.

\begin{table}[H]
\centering
\begin{tabularx}{\textwidth}{@{}l X X@{}}
\toprule
\textbf{Instrument} & \textbf{Rol} & \textbf{Context de utilizare} \\
\midrule
Claude Code (Anthropic) & Asistent de programare interactiv & În terminal, pe parcursul dezvoltării \\
Google Gemini API & Motorul de decizie al agenților & Integrat în produs ca funcționalitate centrală \\
DeepEval + Gemini & Model-judecător pentru evals & În suita de testare a componentei AI \\
\bottomrule
\end{tabularx}
\caption{Instrumentele AI utilizate în proiect}
\end{table}

\subsection{Claude Code --- asistent de dezvoltare}

Claude Code a fost folosit ca \textit{pair-programmer} în terminal, cu acces la
fișierele proiectului. Spre deosebire de autocomplete, operează la nivel de
\textbf{task}, nu de linie: primește o descriere în limbaj natural, citește
fișierele relevante, raționează arhitectural și aplică modificări coordonate pe
mai multe fișiere.

\textbf{Arii susținute:}
\begin{itemize}
  \item \textbf{Arhitectură backend} --- reprezentarea grafică a tablei
        (\texttt{graph.ts}, \texttt{seafarers-graph.ts}), sistemul de trăsături cu
        matrice de conflicte, și extensia Seafarers fără a rupe logica Classic.
        Decizia cheie --- \texttt{SeafarersContext} ca parametru \textit{opțional} ---
        a permis zero cuplare între moduri.
  \item \textbf{Frontend} --- rescrierea \texttt{CatanBoard.tsx} pentru Seafarers
        (canvas extins, hexagoane exterioare, trasee punctate pentru nave).
  \item \textbf{DevOps} --- build-uri Docker multi-stage, pipeline CI/CD, și
        diagnoza incidentului din Secțiunea~\ref{sec:bug3}, unde cauza rădăcină a
        fost \textbf{demonstrată experimental} prin reproducerea configurației
        vechi în contextul real de build Docker.
  \item \textbf{Testare} --- generarea suitei Vitest și a harness-ului de simulare.
\end{itemize}

\subsection{Google Gemini --- AI în produs}

Gemini este apelat la fiecare tură a fiecărui agent prin
\texttt{backend/src/game/ai.ts}, gestionat de \texttt{gemini-client.ts} cu un pool
de 8 chei API în rotație. La eroare recuperabilă (\texttt{429}, \texttt{503},
timeout), clientul rotește automat cheia.

\subsubsection{Sistemul de personalitate}

Fiecare agent are trăsături alese din \textbf{15 posibile}: \textit{Empathic,
Greedy, HardTrader, Aggressive, Expansionist, Defensive, SmartBuilder,
FastSpender, DevFocused, Diplomatic, Opportunist, Militarist, Settler, Merchant,
Tactician}. Trăsăturile sunt stocate relațional și injectate în system prompt,
producând comportamente strategic distincte. Perechile incompatibile sunt
respinse la creare, pe baza tabelei \texttt{TraitConflict}.

\subsubsection{Utilizări suplimentare}

\begin{itemize}
  \item \textbf{Moderare de conținut} --- numele agenților și meciurilor sunt
        validate înainte de stocare, cu o listă locală de termeni interziși ca
        primă barieră (funcționează și când API-ul este indisponibil);
  \item \textbf{Rezumate narative} --- la finalul fiecărui meci se generează un
        rezumat stocat în \texttt{Match.summary};
  \item \textbf{Comentariu live} --- o frază per acțiune, difuzată prin SSE.
\end{itemize}

\textbf{Estimare de consum:} $\approx$40--80 de apeluri per meci pentru decizii,
plus unul pentru rezumatul final. Fallback-ul determinist garantează că jocul
funcționează și fără nicio cheie configurată.

\subsection{Lecții învățate}

\subsubsection{Ce a funcționat bine}
\begin{itemize}
  \item utilizarea asistentului la nivel de task (nu de linie) --- semnificativ
        mai productivă pentru modificări cross-file;
  \item parametrul opțional \texttt{SeafarersContext} --- extindere fără regresie
        pe modul Classic;
  \item fallback-ul determinist --- dezvoltare și testare fără consum de cotă API;
  \item \textbf{verificarea empirică a ipotezelor} --- în incidentul de producție,
        cauza rădăcină nu a fost presupusă, ci reprodusă și demonstrată.
\end{itemize}

\subsubsection{Ce a necesitat judecată umană}
\begin{itemize}
  \item decizii de schemă (când \texttt{--accept-data-loss} este acceptabil față
        de o migrare propriu-zisă);
  \item granițele arhitecturale între modurile de joc;
  \item compromisuri cost/calitate în verbozitatea prompt-urilor.
\end{itemize}

\subsubsection{Direcții viitoare}
\begin{itemize}
  \item trecerea de la \texttt{prisma db push} la migrări versionate
        (\texttt{prisma migrate deploy});
  \item paginarea log-ului de evenimente și transmiterea incrementală prin SSE;
  \item caching pentru stări de joc repetitive, pentru reducerea costului API;
  \item instrumentarea latenței deciziilor AI.
\end{itemize}

% =====================================================================
\section{Concluzii}
% =====================================================================

Proiectul a livrat o platformă funcțională, desfășurată în producție, care
acoperă toate cele 15 user stories din backlog. Din perspectiva
\textit{Metodelor de Dezvoltare Software}, elementele cele mai relevante sunt:

\begin{itemize}
  \item \textbf{Proces} --- backlog structurat pe epice, prioritizare MoSCoW,
        dezvoltare pe feature branches, integrare exclusiv prin Pull Requests
        cu review;
  \item \textbf{Modelare} --- sistemul a fost descris prin șapte diagrame UML
        complementare, acoperind atât structura, cât și comportamentul;
  \item \textbf{Calitate} --- trei niveluri de testare, inclusiv testare bazată pe
        proprietăți pentru simulare și \textit{evals} pentru componenta
        non-deterministă;
  \item \textbf{Automatizare} --- pipeline CI/CD cu șase verificări obligatorii
        înaintea livrării automate în producție;
  \item \textbf{Învățare din incidente} --- căderea de producție analizată în
        Secțiunea~\ref{sec:bug3} a produs o îmbunătățire permanentă de proces,
        nu doar o remediere punctuală.
\end{itemize}

Cea mai importantă lecție tehnică a proiectului este că \textbf{o suită de teste
verde nu garantează un produs funcțional} dacă mediul validat diferă de artefactul
livrat --- observație care a schimbat structura pipeline-ului.

\end{document}
